\chapter{RELATED WORK}

Hole filling on triangle meshes has a long history. The papers that turned out to matter for this project fall into three rough groups: combinatorial methods that triangulate the hole as a polygon, variational methods that solve a global energy on the patch, and normal-field methods that propagate orientation information from the boundary inward. NFD borrows from all three and is closest in spirit to the third.

\section{Combinatorial fillers}

The ear-clipping family is the oldest. Meisters~\cite{meisters1975} proved in 1975 that every simple polygon has at least two ears, which is the topological fact that makes ear clipping work. Held's FIST system~\cite{held2001} is the robust industrial implementation, and most mesh-tool ear clippers descend from it either directly or by reimplementation. The version in this report borrows the layered fallback structure from FIST without using its exact-arithmetic predicates, because the patch sizes encountered here are small enough that ordinary floating-point comparisons do not cause trouble.

For 3D hole filling the canonical reference is Liepa~\cite{liepa2003}, an advancing-front method that picks each new triangle directly in 3D rather than projecting to 2D and lifting back. Liepa is what MeshLab's \texttt{Close Holes} filter implements, and it is the direct baseline this report compares against in Chapter~6. The patch it produces is geometrically flat in the following sense: each new triangle is chosen on local quality criteria (dihedral angle, edge length, balance of the front) without any input from the curvature of the surrounding mesh. On a sphere this is immediately visible.

Barequet and Sharir~\cite{barequet1995} formulate the same problem as a constrained Delaunay triangulation, which is theoretically cleaner but expensive and not native to the MeshLab/VCGlib tooling. The standard approximation is Lawson-flip iteration, which locally improves quality without solving the global problem. The implementation here applies it twice, once after initial ear clipping and once after centroid refinement.

The combinatorial line shares one structural limit. The patch is fully determined by the boundary polygon and the heuristic, so it cannot reflect curvature information from outside the polygon --- there is no mechanism for that to enter. The next two groups of work address this in different ways.

\section{Variational and PDE methods}

The simplest variational approach treats the patch as a discrete minimum surface and solves $\Delta u = 0$ on patch positions with the boundary fixed. It is fast, gives a smooth patch, and biases toward flatness because it minimises area. On a sphere the patch sits inside the chord rather than on the sphere; the dimple problem of ear clipping is reproduced through different machinery. Botsch and Kobbelt's Laplacian deformation framework~\cite{botsch2004}, although the paper itself is about interactive editing rather than hole filling, is built on the same Laplacian-with-Dirichlet machinery and shares the same shrinkage tendency when adapted to hole filling.

Going to a fourth-order energy --- bi-Laplacian, equivalent to thin-plate splines --- lets the patch encode boundary tangents in addition to positions, and produces $C^2$ patches that can bulge outward when the surrounding surface is curved. Nealen et al.~\cite{nealen2006} present a unified Laplacian/bi-Laplacian framework that covers this as one of several modes. The cost is a denser, worse-conditioned linear system: the bandwidth of $\Delta^2$ is roughly the square of $\Delta$ on a triangle mesh, and direct solvers slow down accordingly. For a plugin that runs inside MeshLab's interactive loop this lives on the wrong side of the speed/quality trade-off.

Of all the references in this section, the closest in spirit to NFD is Pernot, Moraru, and V{\'e}ron~\cite{pernot2006}. Instead of minimising a quadratic energy on positions, they minimise the variation of curvature across the patch, simulated as a thin-plate mechanical model whose boundary forces are driven by the surrounding curvature. The high-level intuition --- carry curvature information from the boundary inward --- is the same as NFD's, but the formulation is on positions, the iteration is non-linear, and the input must be smoothly parameterisable. Reading their paper alongside Verdera et al.\ helped clarify what is actually being asked of the patch: not just a smooth surface, but one whose curvature continues that of the boundary.

Poisson surface reconstruction~\cite{kazhdan2006, kazhdan2013} occupies a different niche. It is not strictly a hole filler --- the input is a point cloud or oriented mesh and the output is a globally smooth surface, with holes filled as a side effect. For hole-repair purposes the difficulty is exactly that. Parts of the input that did not need fixing get reshaped along with the broken parts, and the input-output triangle correspondence is lost in the volumetric intermediate.

Laplacian Surface Editing by Sorkine et al.~\cite{sorkine2004} is mentioned here because of a useful conceptual distinction it sharpens: the Laplacian can be used as an \emph{operator} on positions or as a \emph{representation} of geometry through differential coordinates. Minimum-surface and bi-Laplacian hole fillers use the Laplacian as an operator. Differential-coordinates methods use it as a representation. NFD belongs to the first camp, although its operator is applied to normals rather than positions.

\section{Normal-field and curvature-based methods}

Verdera, Caselles, Bertalm{\'i}o, and Sapiro~\cite{verdera2003} is the closest precursor and where the core idea of NFD originates. They diffuse each component of the surface normal field across the hole through a non-linear PDE, with boundary normals as Dirichlet data, and reconstruct positions afterwards. NFD is essentially this idea translated from level-set surfaces (where they discretise the PDE on a regular volumetric grid) to triangle meshes (where the cotangent Laplacian gives the right discrete operator directly). The level-set framework was a reasonable choice in 2003 because the triangle-mesh discrete-differential-geometry tooling was still settling: Pinkall--Polthier was known but Meyer et al.\ had only just appeared, and not every implementation used cotangent weights as a default. Translating to triangle meshes today is more or less mechanical.

The diffusion idea itself comes from earlier work on image inpainting by Bertalm{\'i}o, Sapiro, Caselles, and Ballester~\cite{bertalmio2000}, from the same research group. Verdera et al.\ is the surface-normals analogue of Bertalm{\'i}o et al.'s pixel-intensities approach. Knowing this connection was useful while reading the surface-holes paper, because the level-set notation is dense and the underlying intuition is clearer when read alongside the 2D image case.

Local normal-interpolation methods are a parallel branch that interpolates boundary normals into the interior without a global PDE. Wang and Oliveira~\cite{wang2007} use radial basis functions, Jun~\cite{jun2005} uses a moving-least-squares approach with sub-region partitioning. Without a global smoothness criterion these methods are simpler to implement but degrade on larger or non-convex holes; the seams between sub-regions in Jun's piecewise approach in particular tend to need separate cleanup.

Zhao, Gao, and Lin~\cite{zhao2007} sit between Verdera et al.\ and the local methods. They use a Poisson-style PDE for fairing patch positions, combined with a separate boundary-driven displacement step. The two-stage structure is the same as Stages~4 and~6 in this report, but the variable they solve for is position, not normal, so the formulation inherits the shrinkage of position-based methods --- partially compensated by the displacement step but not eliminated.

\section{Foundations}

The discrete operators NFD relies on come from the broader discrete-differential-geometry literature. The cotangent Laplacian was introduced by Pinkall and Polthier~\cite{pinkall1993} in 1993 for discrete minimal-surface construction, and was packaged as a standard geometry-processing operator by Meyer et al.~\cite{meyer2003} a decade later, with the mean-curvature normal identity that Stage~2 of NFD uses. Wardetzky et al.~\cite{wardetzky2007discrete} gave the useful negative result that no discrete Laplacian satisfies all desirable properties simultaneously, which retroactively explains why several different cotangent-style operators coexist in the literature. The version used here is convergent and intrinsic but breaks the maximum principle on obtuse triangles, and NFD handles this by clamping negative weights to zero. Sharp and Crane~\cite{sharp2020} address the same problem more cleanly through intrinsic Delaunay remeshing; that would be a worthwhile upgrade for inputs with many obtuse triangles.

Two other operators are reused as-is. Desbrun et al.~\cite{desbrun1999} introduced implicit fairing by diffusion in 1999, which is essentially the numerical recipe NFD uses in Stage~4 (assemble the cotangent Laplacian, solve the implicit-Euler system, iterate) but applied to vertex positions rather than normals. Taubin's $(\lambda, \mu)$ filter~\cite{taubin1995} is the standard non-shrinking low-pass smoother, used in Stage~6 of NFD because uniform Laplacian smoothing would flatten the dome. Crane et al.~\cite{crane2013heat} proposed the heat method for geodesic distance computation, which would be a more accurate replacement for the multi-source Dijkstra in Stage~6; we have not measured how much accuracy it would buy in practice on the test patches.

The mesh-repair survey of Attene, Campen, and Kobbelt~\cite{attene2013} is the right reference for context. Hole filling is one of many mesh-defect classes; a production-grade repair tool needs to handle the others (non-manifold edges, self-intersections, inverted normals) as well. NFD assumes the input has been cleaned of those defects, which is reasonable inside MeshLab where standard cleanup filters are one menu away.

\section{Where NFD sits}

The components of NFD are individually well-known: heat-equation diffusion of normals is from Verdera et al., the cotangent Laplacian is standard, Taubin smoothing dates from 1995. What this report contributes is putting them together in a way that runs end-to-end inside MeshLab without voxelisation, with an implicit-Euler solver and a one-time $LDL^T$ factorisation that keeps the inner loop fast, a spherical-cap displacement model whose dome height is derived from boundary geometry rather than supplied as a parameter, and a $\sqrt{2t-t^2}$ profile that gives tangent-plane continuity at the patch boundary --- the property the more common parabolic profile fails to provide.

Among the closest references --- Verdera et al., Pernot et al., Zhao et al.\ --- none combine these specific choices in the same way. The result is best understood as a practical translation of Verdera-style normal diffusion into the triangle-mesh setting, picked apart into a sequence of stages that each have a clear role and a known fallback when they fail.
